\chapter{統合効果の要因分析と特徴利用の考察}\label{ch:discussion}
第4章では，全域回帰と総見逃しへの利得が，ONB近傍や検知到達段階の一律改善にはつながらないことを示した．本章では，まず保存予測から直接確認できる誤差の変化を示し，次に入力依存との関係を考察する．最後に，モデル上の依存から物理的な音源へ説明を進めるために何が不足するかを明らかにする．保存予測の分析は10月15日の主実行と10月14日の接続分析\cite{localLink}，入力対照は10月7日の完了記録\cite{localInput}に基づく．

\section{熱流束領域別の予測誤差と偏り}
\subsection{全域改善が生じた領域}
表\ref{tab:regions}では，評価540区間をONB前225，ONB近傍30，近傍を除く高熱流束側285区間の非重複な3領域へ分ける．無雑音では元3から5への変更により，近傍と高熱流束側のRMSEが小さくなる一方，ONB前は27.10から29.54\,\kw{}へ増えた．$-20$\,dBではONB前と高熱流束側が改善し，近傍は93.92から99.32\,\kw{}へ悪化する．全域改善の内訳は雑音条件で異なる．

\begin{table}[htbp]
\centering\small
\caption{非重複な領域別の元3 MSEから5モデルへのRMSE変化 [\kw{}]．}
\label{tab:regions}
\begin{tabular}{lrrrrr}
\toprule
領域 & $n$ & 無雑音・元3 & 無雑音・5 & $-20$・元3 & $-20$・5 \\
\midrule
全域 & 540 & 34.07 & 27.80 & 81.14 & 75.23 \\
ONB前 & 225 & 27.10 & 29.54 & 102.01 & 94.29 \\
ONB近傍 & 30 & 39.58 & 16.45 & 93.92 & 99.32 \\
高熱流束側 & 285 & 38.14 & 27.33 & 57.70 & 51.63 \\
\bottomrule
\end{tabular}
\source{領域分けを明示した保存予測分析\cite{localLink}．ONB以上全315区間の指標と混同しない．}
\end{table}

非重複な領域$r$の標本数を$n_r$とすると，$\mathrm{MSE}_{\rm all}=\sum_r(n_r/N)\mathrm{MSE}_r$である．RMSE自体の加算ではない．近傍は全体の30/540であり，その領域が悪化しても，他の領域の二乗誤差減少が上回れば全域値は改善する．これは近傍を軽視してよいという意味ではなく，全域目的と検知目的の違いを理解するための関係である．

\subsection{近傍の回復と判定が一致しない理由}
主実行では，$-12$から$-20$\,dBへ進むと5モデルの近傍biasは$-107.52$から$-94.79$\,\kw{}へ変わる．過小予測が弱まることで近傍RMSEは小さくなるが，予測がONB閾値に届かなければFNは残る．さらに，他のONB以上の段階でFNが増えると総見逃しは悪化する．局所RMSEと判定件数が異なる動きをすることは，指標の計算誤りを意味しない．

保存重みでモデルごとの平均予測変化を分けると，この近傍biasの上方変化12.72\,\kw{}に対し，ETの重み付き寄与は約8.02\,\kw{}である\cite{localMainRun}．これは固定予測の代数的な内訳であり，ETが気泡音を適切に取り出したことや，雑音が物理的に検知を助けたことの証明ではない．

全域MSEで配分を決めていることと，近傍性能が一律には良くならない観測は整合する．ただし，全域目的だけが近傍の悪化を生んだと因果的に確定するには，入力・学習器・分割等を揃えて目的関数だけを変える比較が必要である．本稿では，観測事実とそれに整合する解釈の段階で述べる．

\section{モデル間の予測誤差の関係と補完性}
\subsection{逆符号の相殺と同方向誤差の縮小}
平均化の説明として，過大予測と過小予測の相殺が挙げられる．しかし，現在の統合効果はそれだけでは説明できない．元3・HGB・ETの3予測群の誤差方向で分けた結果を表\ref{tab:directions}に示す．無雑音では同方向群が全体MSE減少の54.65\%，$-20$\,dBでは87.15\%を占めた．
\begin{table}[htbp]
\centering\small
\caption{元3・HGB・ETの3予測群による誤差方向別の全体MSE減少への寄与．MSE寄与の単位は\mseunit{}．}
\label{tab:directions}
% Generated from saved evidence; do not hand-edit.
\begin{tabular}{llrrrr}
\toprule
条件 & 誤差群 & $n$ & 改善数 & MSE減少寄与 & 割合 [\%] \\
\midrule
無雑音 & 逆符号を含む & 211 & 178 & 175.88 & 45.35 \\
無雑音 & 同方向 & 329 & 204 & 211.97 & 54.65 \\
$-20$ dB & 逆符号を含む & 112 & 86 & 118.61 & 12.85 \\
$-20$ dB & 同方向 & 428 & 271 & 804.44 & 87.15 \\
\bottomrule
\end{tabular}

\end{table}

例えば，三つの予測が全て過小でも，その中に正解に近い予測があれば，統合によって基準の大きな過小誤差を縮められる．正解を挟む予測がなければ誤差を完全に相殺できないが，基準モデルより良くなることはある．同方向群の改善を「相殺が働いた」と一括して呼ぶと，この違いが失われる．

表の87.15\%は，同方向群の二乗誤差減少の総和を全体の二乗誤差減少で除した割合である．群内の87.15\%が改善したという意味ではなく，音響特徴の寄与率でもない．また，五つの単体全てが同方向という分類ではなく，実際の追加統合単位である三つの予測群の分類である．

\subsection{調整方向と大きさの分解}
式\eqref{eq:decomp}による無雑音のMSE差は，
\begin{equation}
 -851.94+464.09=-387.85\quad[\mseunit{}]
\end{equation}
であり，$-20$\,dBでは，
\begin{equation}
 -1346.86+423.81=-923.05\quad[\mseunit{}]
\end{equation}
となる\cite{localLink}．元誤差と調整の積による負の項が，調整自体の二乗による正の項を上回ったことが，全体改善の直接の説明となる．この分解は保存予測に対して恒等的に成立するため，性能差をどの調整が作ったかを追跡できる．一方，なぜその予測をモデルが出したかという入力・物理の説明は，別の証拠を要する．

\section{統合による誤りの改善と残存要因}
\subsection{訂正された判定と失われた判定}
元3 MSEから5への変更では総FNが減るが，ET単体から5への変更では陽性が平均化で失われる例がある．この違いは，全域回帰に有利な出力調整と，閾値通過を保つ調整が同じではないことを表す．第4章の無雑音例ではETの2 FNを訂正し，10件を新たにFNとした．良い単体の判定をそのまま保持する規則として，現在の凸結合を解釈してはならない．

5モデルで残る無雑音27 FNのうち，11件は五つの単体全てが陰性であり，16件には少なくとも一つ陽性の単体があった\cite{localMainRun}．前者は現在の単体群に共通した失敗，後者は単体の陽性が統合で残らなかった失敗である．ただし，この診断は正解を見た後の分類であり，16件を運用時に自動的に見分けて全て救えることを示さない．

全単体共通の失敗も，入力に情報がないことの証明ではない．現在の表現，学習器，損失，学習標本では利用できなかった可能性が残る．入力除去後の別モデルで一部が訂正される例もあるため，「五つが失敗したから音響に原理的な限界がある」と断定しない．

\subsection{4モデルから5モデルへの追加効果}
元3＋ETから5への変更は，保存された配分において，概ねHGBとETの予測差の27.25\%を加える調整である．10月14日の保存予測では，無雑音のMSE差が$+24.94$\,\mseunit{}，$-20$\,dBが$-132.91$\,\mseunit{}となった\cite{localLink}．前者は追加調整が誤差を大きくし，後者は誤差を減らす方向へ働く．

$-20$\,dBでは改善区間268に対して悪化区間272であっても，全体MSEは小さくなった．件数の多数決ではなく，各区間でどれだけ誤差が変わったかが二乗誤差を決めるためである．「改善した標本が多い」「モデルが一つ増えた」といった説明だけでは，追加効果の量は分からない．

\section{入力特徴に対するモデルの依存性}
\subsection{1秒内の時間配置}
表\ref{tab:shuffle}に，無雑音の共通時間行並べ替えの結果を示す．CNN＋TransformerのRMSEは35.34から47.99\,\kw{}，AlexNetは39.88から56.45\,\kw{}へ増えた．HGBとETは不変である．元3 MSEは34.07から47.52\,\kw{}，5モデルは27.80から31.10\,\kw{}となった．
\begin{table}[htbp]
\centering\small
\caption{固定モデルへの共通時間行並べ替え．全域RMSE [\kw{}]．元の10月7日対照を用いる．}
\label{tab:shuffle}
% Generated from saved evidence; do not hand-edit.
\begin{tabular}{lrrr}
\toprule
方式 & 元入力 & 時間行並べ替え後 & 増加量 \\
\midrule
RF (PCA+XGBRF) & 71.71 & 71.21 & -0.50 \\
CNN+Transformer & 35.34 & 47.99 & 12.65 \\
AlexNet & 39.88 & 56.45 & 16.58 \\
HGB & 33.25 & 33.25 & 0.00 \\
ET & 28.40 & 28.40 & 0.00 \\
5-model & 27.80 & 31.10 & 3.30 \\
3-MSE & 34.07 & 47.52 & 13.45 \\
\bottomrule
\end{tabular}

\end{table}

HGB/ETの不変性は，34特徴が時間行の順番を保持しない構造と一致する．CNN系の変化は，固定モデルの予測が1秒内の配置に依存することを支持する．ただし，並べ替えは連続した音響として不自然な配置を作り得る．その誤差増加を，実際の気泡成長・凝縮過程の情報を正しく使っている量としては解釈できない．

5モデルの加工後の悪化が元3より小さいことは，順序不変の予測群を大きく含む固定配分と整合する．しかし，並べ替えと水流音付加は異なる加工である．この結果から，時間順序への依存差が実際の水流雑音下の利得を生んだ因果割合まで確定することはできない．

\subsection{強度と帯域への感度}
全域powerを2倍にした対照では，HGBのRMSEは33.25から71.33\,\kw{}，ETは28.40から76.58\,\kw{}へ増えた．2.1--2.5\,kHzのpowerを0.1倍にした対照では，それぞれ147.97，133.15\,\kw{}となった\cite{localLink}．順序に不変であることは，強度や帯域の変化にも不変であることを意味しない．

この帯域には絶対powerと相対形状の複数の特徴が関係する．加工による悪化は，元モデルがその情報に依存していたことを示すが，その帯域だけが沸騰の情報を持つという証拠ではない．全域gainと帯域加工が特徴分布を変える点，帯域座標がresize後の近似である点にも留意する．物理周波数の精密な同定には，元波形・伝達特性との照合が必要である．

\section{入力表現・統合方法と予測特性の関係}
\subsection{絶対強度と相対形状を除去した対照}
34特徴を絶対強度5特徴または相対形状等29特徴へ減らして再学習した結果を表\ref{tab:removal}に示す．両モデルとも全34特徴が，全域回帰で削除表現より小さい誤差を与えた．この範囲では，強度と形状を併用する表現が回帰精度を支える．
\begin{table}[htbp]
\centering\small
\caption{特徴除去後に再学習した単体の端点性能．RMSE [\kw{}]，FN/FPは区間数．}
\label{tab:removal}
\begin{tabular}{llrrrr}
\toprule
モデル & 入力 & 無雑音RMSE & 無雑音FN/FP & $-20$ RMSE & $-20$ FN/FP \\
\midrule
HGB & 全34 & 33.25 & 21/0 & 74.19 & 49/0 \\
HGB & 絶対5 & 44.71 & 25/2 & 129.78 & 30/106 \\
HGB & 相対等29 & 46.46 & 18/0 & 93.77 & 52/0 \\
ET & 全34 & 28.40 & 19/2 & 75.43 & 50/1 \\
ET & 絶対5 & 40.85 & 29/1 & 120.38 & 24/113 \\
ET & 相対等29 & 41.77 & 23/1 & 80.09 & 51/0 \\
\bottomrule
\end{tabular}
\source{入力除去後の保存対照\cite{localInput,localLink}．}
\end{table}

絶対5のHGBは$-20$\,dBで全34よりFNが少ない一方，FPが106へ増える．誤差やFN一つだけで「情報を減らした方が検知に有効」と結論できない例である．また，相対形状等29には時間power変動係数も含まれるため，完全に周波数形状だけの対照と呼ばない．

\subsection{特徴変更と配分変更の区別}
HGBだけを相対等29へ変更した5統合では，元の重みを固定すると無雑音RMSEが27.80から29.06\,\kw{}へ増え，6月18日の$q_{100}$は376.32\,\kw{}のままであった．一方，OOFで重みを再fitするとRMSEは27.39\,\kw{}，$q_{100}$は322.11\,\kw{}となる．このときHGBの配分は約27.25\%から5\%へ下がり，ETへ70\%が配分された\cite{localInput}．

従って，この$q_{100}$の前進を「相対形状を使ったため」と単独で説明することはできない．入力変更によるclean MSE差は元重みで$+71.78$\,\mseunit{}，そこから重み変更で$-94.56$\,\mseunit{}であり，配分変更を伴う結果である．元3＋ETにも既に同じ$q_{100}$の利点があるため，新表現だけの独自な検知利得とも断定しない．

元3内比率を自由にする既存対照では，無雑音RMSEは27.47\,\kw{}となるが，$-20$\,dBでFNは52から53へ増える．固定比率の解除にも交換関係がある．現在の構成は，全域回帰，判定，元3保持という比較意図を合わせた採用条件であり，全ての指標の普遍的最適解ではない．

\section{特徴利用の物理的解釈とその限界}
\subsection{モデル依存と物理音源をつなぐ条件}
ここまでの結果から直接言えるのは，モデル群が同じ入力変化に同じ応答をするわけではなく，統合による予測調整が領域ごとの誤差を変えることである．一方，帯域への感度が高いことだけでは，その帯域が気泡の成長・凝縮あるいは固体振動を表すと確定できない．原音の発生，伝達，収録，前処理，学習という複数段階が観測特徴へ関係するためである．

手塚・植木らの2024年発表要旨\cite{tezuka2024}には，細線径に伴う音響ピークの移動と加熱固体の共振帯との対応が記述されている．一方，手塚の2026年最終修士論文\cite{tezuka2026}の第5章では，調査範囲で共振と重なる部分はあるが伝熱構造体振動との強い相関は確認できなかったと述べ，2\,kHz付近のピークを気泡凝縮に伴う圧力変動として考察している．対象と証拠の追加により解釈が変わり得るため，初期要旨だけを使って本研究の帯域を固体共振と断定しない．本データの音源を同定するには，装置・計測・前処理条件を照合し，その解釈に対応する独立した根拠を確認する必要がある．

Dunlapら\cite{dunlap2023}では低周波特徴の重要性が報告される一方，本研究では500\,Hz高域通過後の入力を用いる．異なる対象・装置・前処理の結果であり，一方が他方を直ちに否定する関係ではない．本研究のモデルが3\,kHz以下を利用した事実を，低周波域が物理的に不要であるという結論へ広げない．

\need{L02}{手塚の最終修論の図3.2.1，第3章のサブクール度比較，第4章の同期観測を，対象データの原条件と対応づけて検討する．必要に応じて引用元の査読論文まで確認する．現在の2.1--2.5\,kHz帯を特定の音源へ確定できる根拠は不足している．}

\subsection{説明性の到達段階}
説明性は，少なくとも三段階に分けると現在地を明確にできる．第一は，保存予測のどの調整が誤差を変えたかという出力上の説明であり，式\eqref{eq:decomp}と誤り群の分解で直接確認できる．第二は，どの入力加工が出力や誤差を変えるかというモデル上の依存であり，共通加工と特徴除去の対照がこれを支える．第三は，その特徴がどの沸騰機構を捉えたかという物理的な説明であり，独立した観測と実験条件の照合を要する．

IG\cite{sundararajan2017}を用いる場合にも，出力が回帰値であること，基準入力，積分近似の収束と寄与和の整合を確認する必要がある．数値的に収束した寄与画像であっても，それだけで物理音源を同定したことにはならない．本稿の中心結論は共通入力対照と保存予測で支え，未確認のIG画像の印象によって補強しない．IGの位置づけと残る確認は付録に示す．

\subsection{現在の説明と追加検証を要する主張}
本研究が現在支持する説明は，異なる依存を持つモデル群の予測を統合することで，逆符号の相殺と同方向誤差の縮小が生じ，一部の領域・判定を改善した，という範囲である．入力依存の差が水流音下の利得を生む因果割合，特定帯域の音源，実際の気泡発生秒に対する応答は未確定である．

\need{V02}{未知WAV・未知日への一般化や統計的な優位を主張する場合は，分割単位と開発への参照履歴を定めた独立評価が必要である．540区間を独立した540実験と見なす推定や，seedの反復だけで未知条件への一般性を代替することはできない．}
\need{V03}{時間配置への依存差が水流雑音時の利得の原因である，または特定帯域が特定の物理音源であると強く結論する場合は，表現・学習器を統制した比較や独立観測を要する．本稿の限定した結論を記述するために，これらを一律の追加実験として要求しない．}

本章の考察により，何が改善したかに対しては出力上の説明とモデル上の依存差を示せるが，物理的な確定には不足が残ることを明らかにした．この区別を保つことで，第1章の目的に対する到達点と未達を，結言で過不足なく述べることができる．
